\section{Discussion}

The results illustrate how repeated unsupervised learning can describe financial relationships within the JSE Top 40 and show where those relationships change. The UMAP hierarchical model produced separated groups in each reduced representation, while the feature profiles and membership comparisons provided a more detailed account of their financial characteristics. The three-date design adds information that was not available in the earlier single-window analysis: whether a company remains with the same peer group, joins a different group or returns to an earlier group.

The model comparison also shows why internal validation must be interpreted carefully. UMAP produced the higher embedding Silhouette Score at every date, but t-SNE produced higher financial-feature scores in June and December 2025. A visually clear map therefore does not establish an equally strong grouping in the original financial variables. The practical value of the representation lies in making relationships easier to examine, while the financial-feature scores indicate how strongly those relationships are supported by the measured inputs.

For equity analysis, the membership histories should be used as a screening layer. ABSA and Old Mutual return to an earlier aligned group, while Reinet and Naspers follow different paths. A change in cluster membership does not prove that a company has become mispriced or experienced an economic structural break. It identifies a relationship that deserves further investigation. An analyst can then examine changes in return, volatility, valuation, scale and trading activity, together with changes in the company's peers.

Interpretability is necessary because a cluster label alone provides limited insight. The median profiles show which financial variables distinguish a group, while company-level histories show how these attributes change over time. The surrogate SHAP analysis provides an additional explanation, but its usefulness depends on fidelity to the cluster labels. Mahalanobis distance gives a separate view of firms far from the multivariate centre. These components answer different questions and make the screening output easier to examine without treating agreement between them as automatic proof of an anomaly.

The South African contribution is a repeated, interpretable financial-clustering case study in a compact equity universe with changing index membership, dual-listed securities and substantial international exposure. Effective-dated membership distinguishes index reconstitution from peer-group movement, while the corrections to share units and Investec line counts show why financial definitions matter in longitudinal comparisons. The method provides a structured way to follow financial similarities over time, with its conclusions conditional on the inputs and modelling choices used.
